\documentclass[12pt]{article}

\usepackage{sbc-template}
\usepackage{graphicx,url}
\usepackage[utf8]{inputenc}
\usepackage[english]{babel}
\usepackage{multirow}
\usepackage[T1]{fontenc}
%\usepackage[latin1]{inputenc}

     
\sloppy

\title{K-ORCH: an orchestrator for distributed multikernel instances}

\author{Eduardo Kaluf\inst{1}, Carlos Alberto Maziero\inst{2}}

\address{
  Departamento de Informática -- Universidade Federal do Paraná (UFPR)\\
  Curitiba -- PR -- Brazil
  \email{eduardo@kaluf.com.br, maziero@inf.ufpr.br}
}

\begin{document}

\maketitle

\begin{abstract}
  The multikernel architecture runs multiple independent kernel instances, each on a dedicated partition of cores and memory of
  the same processor, being a middle ground between containers and VMs and offering predictable
  performance, kernels with specialized configurations, and performance isolation between instances, without CPU and
  memory virtualization overhead.
  Current implementations of this architecture \cite{wang2025multikernel}
  manage these instances only within a single machine, and no tools exist to orchestrate them in distributed environments.
  This work proposes K-ORCH, an orchestrator that provisions, places, and monitors multikernel instances across
  multiple nodes, based on a desired-state model, agents responsible for partitioning cores and memory on demand,
  and placement policies driven by the available physical resources.
  K-ORCH will be experimentally evaluated through a quantitative analysis of provisioning time, control-plane overhead,
  density, and performance interference between instances, compared to orchestrated microVMs, in addition to comparing
  the orchestrator's own static and dynamic partitioning models.
  The expected outcome is to demonstrate the feasibility of distributed multikernel orchestration and to highlight its
  costs, advantages, and limits relative to current alternatives.
\end{abstract}

\section{Introduction} \label{sec:introducao}
% Initial context, multiple cores and motivation for the creation of the multikernel architecture (history)
The number of cores per processor grows incessantly, and processors already exceed a hundred cores on a single chip~\cite{amd2024epyc}.
Monolithic operating systems, designed around data and control structures shared by all cores, have kept pace with this growth
at the cost of increasingly complex mechanisms, such as fine-grained locks, per-core data structures, and CPU affinity and
isolation policies~\cite{boydwickizer2010mosbench}.
These solutions increase complexity and the kernel surface, and make state sharing between cores a bottleneck for
performance and scalability~\cite{boydwickizer2008corey}.
In this context, the multikernel architecture emerges, proposing to treat the machine as a network of independent cores,
with replicated state and explicit message-based communication, instead of a single kernel with shared
state, bringing the kernel itself closer to the idea of a distributed system~\cite{baumann2009multikernel,wentzlaff2009fos}.

% Defining the types of multikernel
Since then, two ways of logically subdividing the cores have been widely explored.
The first restructures the OS entirely, as in Barrelfish, which runs one kernel per core~\cite{baumann2009multikernel}.
The second line, adopted by this work, partitions the physical machine
among multiple independent kernel instances, each with its own cores and memory region,
with the possibility of having different configurations and fine-tuning.
This line emerged with high-performance computing (HPC) in mind, with systems that run lightweight kernels alongside Linux~\cite{shimosawa2014ihk,park2012fusedos,gerofi2016survey},
and was recently proposed for Linux itself, which now loads independent instances on reserved cores and memory, with run-time
reallocation~\cite{wang2025multikernel,wang2025lpc}.
In this model, a host kernel partitions the resources and starts the
instances, which communicate through explicit message passing and shared memory regions, in addition, they access the
rest of the hardware, in general, through virtual devices provided by the host kernel.
In this work, \textit{multikernel instance} refers to this set of objects: dedicated cores, its own memory region,
an application image, and its own kernel configuration.

% Multikernel as a middle ground between containers and virtual machines
Furthermore, the multikernel architecture occupies a unique point between the two forms of resource subdivision that are widely used today.
Containers are executable units of code that package an application and its dependencies, using mechanisms of the OS itself, such as \textit{namespaces} and \textit{cgroups},
to divide resources and limit them.
Containers are lightweight, but they share the host kernel, which widens the attack surface between
applications and leaves them subject to interference from neighbors (noisy neighbors).
On the other hand, VMs run their own
kernel on top of a hypervisor and offer stronger isolation, at the cost of an overhead from guest-host communication,
I/O, and memory access, that is, they have higher resource consumption per instance.
Thus, a multikernel instance comes in as a third way.
It also runs its own kernel, but directly on the physical cores, without the CPU and memory virtualization layer of VMs, which favors
predictable performance and allows the kernel to be specialized for specific workloads.
On the other hand, allocation has whole-core granularity, and security isolation between instances depends on the
hardware mechanisms of the machine.

% Point out distributed systems and orchestration as the main adoption factor, reinforce that multikernel is restricted to a single machine
The large-scale adoption of containers and VMs is strongly tied to the creation of orchestrators that manage these
units, with emphasis on the line that went from Borg to
Kubernetes~\cite{burns2016borg}.
Existing tools for multikernel, on the other hand, such as IHK~\cite{shimosawa2014ihk} and kerf~\cite{multikernel2025kerf},
offer mechanisms to create, configure, monitor, and destroy instances, similar to container tools such as \textit{docker compose},
but they operate on a single machine, without defining policies about where and when these instances should be created.
Without a robust orchestration layer, the architecture tends to remain restricted to single-machine experiments, far from
distributed production environments such as the cloud, clusters, fog computing, and edge computing.

% K-ORCH proposal transitioning to the topic
Therefore, this work proposes K-ORCH, an orchestrator for multikernel instances on nodes of a local network,
under a single administrative domain, capable of meeting the specific needs of this emerging architecture.
K-ORCH receives requests for
instances described with the desired state, chooses a target node based on available cores and memory, delegates to
an agent on each node the partitioning of resources and the initialization of the instance,
and also monitors the instances in order to detect and recover from failures.
The proposal will be evaluated in terms of
provisioning time, control plane cost, density, and performance interference between instances,
compared with orchestrated microVMs.

% Paragraph on the organization and presentation of the document
The remainder of this document is organized as follows. Section~\ref{sec:just} presents the justification for the
work along with the research question, the general objectives are in Section~\ref{sec:objetivos} and the specific ones in Section~\ref{subsec:objetivo-especificos}.
Section~\ref{sec:related} discusses related work, and Section~\ref{sec:methods} describes the methodology.
Finally, Section~\ref{sec:results} presents the expected results, and Section~\ref{sec:crono}, the two-year
schedule for carrying out the work.

\section{Justification and Motivation} \label{sec:just}

% Evidence that the problem is in fact real
The trade-off between isolation and performance described in Section~\ref{sec:introducao} has concrete consequences in
multi-tenant environments.
Container escape vulnerabilities, such as CVE-2019-5736 and CVE-2024-21626 \cite{cve20195736,cve202421626},
make evident the risk of sharing the same kernel among workloads from different customers.
In addition, interference between workloads leads providers to keep their servers underutilized, in order to
preserve the latency of critical services~\cite{lo2015heracles}.
Industry and academia seek intermediate solutions
to deal with these problems, such as Firecracker~\cite{agache2020firecracker}, developed to run functions and
containers from different customers in microVMs, and Kata Containers~\cite{katacontainers}, which runs each pod (basic unit of Kubernetes)
in a lightweight virtual machine.
These efforts indicate that there is demand for alternatives to the two traditional models, where the multikernel architecture has a lot of
room to establish itself as a third option.

% For whom, scenarios where only multikernel fits
Multikernel is a huge advantage for large-scale applications in which specific OS configurations
directly influence the performance of the instances being run, such as databases and monitors.
This makes it unnecessary to dedicate an entire machine specifically to the applications, something that is not possible
with containers due to sharing of the host kernel, avoiding costs and giving greater speed to the
computation existing in these systems.
In HPC environments, the use of lightweight kernels alongside Linux has already proven to be a viable alternative in massive production systems~\cite{gerofi2016mckernel}.
In edge computing, where hardware is limited and latency must be predictable, execution without the CPU and memory virtualization layer is also
a great advantage.
Both are scenarios in which the use of multikernel instances, with specific configurations, would benefit providers.

The proposal for multikernel support in Linux~\cite{wang2025multikernel} brings the architecture closer
to the most widely used kernel on servers, but without tools that allow it to be expanded beyond a single machine.
Just as the adoption of containers was extremely influenced by orchestrators, the adoption of multikernel in distributed environments depends on an
equivalent layer.
More than that, from a scientific point of view, characterizing the costs and limits of this orchestration is a contribution
in itself, regardless of how the architecture behaves on the proposed criteria.

An alternative to building a dedicated orchestrator would be to integrate multikernel instances into Kubernetes through
an API compatible with its interfaces, similar to what Kata Containers does with its lightweight VMs.
This work chooses to explore a dedicated orchestrator mainly because placement depends on physical resources, such as free cores and
reserved memory.
Therefore, the intention is to investigate partitioning without being tied to an existing platform; in addition, an independent prototype makes it possible
to measure the cost of orchestration in isolation, with less interference and greater accuracy.

% What mechanisms are needed to provision, place, and monitor
% multikernel instances across multiple nodes, and what is the cost relative to orchestrated VMs/microVMs (feasibility + cost/overhead + preserved isolation)?
Therefore, the central question that serves as the basis of this work is: what mechanisms are needed to provision, place, and monitor multikernel instances
across multiple nodes, and what is the cost of this management relative to the orchestration of microVMs?

\section{Objectives} \label{sec:objetivos}

% General: design, implement, and evaluate K-ORCH (answering the research question)
The general objective of this work is to design, implement, and evaluate K-ORCH, which must be capable of provisioning,
placing, and monitoring multikernel instances across multiple nodes.

\subsection{Specific objectives}\label{subsec:objetivo-especificos}

\begin{itemize}
  \item Survey the state of the art: identify the mechanisms of current orchestrators and of multikernel management
  tools, connecting the two areas, in order to determine which mechanisms can be reused and which
  need to be created or adapted specifically for K-ORCH.
  \item Define the orchestrable unit: characterize the multikernel implementation that will be used, including how
  cores and memory are reserved and reallocated, and formalize what makes up an instance and the desired-state model
  that the orchestrator must maintain.
  \item Design the architecture: define the operation of the control plane and its interaction with the data plane,
  design the agent responsible for instantiating the multikernel units from the host kernel of each node,
  and define the communication model between the agents and the orchestrator.
  \item Implement the lifecycle and placement of instances: implement the creation, initialization, stopping, and
  removal of instances and the placement mechanism, responsible for choosing the node that will serve
  each request based on free cores and available memory.
  \item Develop the resource partitioning model: implement a hybrid model that allows choosing between reserving memory and CPU regions at
  the initialization of each node or partitioning dynamically on each request.
  \item Implement instance monitoring: collect metrics on the state and resource usage of each node,
  detect failures and loss of communication with the agent, using monitoring to recover instances with errors,
  relocating or restarting them as necessary.
  \item Experimentally evaluate the orchestrator: measure provisioning time, the overhead of the agent and the
  control plane, instance density per node, and performance interference between instances, comparing the
  results with orchestrated microVMs, and compare static and dynamic partitioning models.
\end{itemize}

K-ORCH will be designed and evaluated for nodes connected on the same local network and under a single administrative domain.
In this scenario, it is assumed that the nodes are administered by the same entity and that the network offers low latency.
The extension to multiple local networks under the same administrative domain, such as in a datacenter, is considered
a natural evolution of the work.
The orchestration of nodes distributed over the Internet, under multiple administrative domains, is outside the
scope of this work.

\section{Related Work} \label{sec:related}

This section presents related work, showing the evolution of the multikernel architecture and connecting other
works to the focus of this proposal: provisioning, placing, and monitoring multikernel instances across different nodes.
Section~\ref{sec:rel-multikernel} presents the evolution of architectures with multiple kernels, from systems for
shared-memory multiprocessors to the multikernels used in HPC environments and the recent proposal for
Linux.
Section~\ref{sec:rel-fundamentos} discusses how these systems decide on resource allocation, how they
monitor them, and how they ensure the isolation and availability of instances.
Section~\ref{sec:rel-orquestracao} addresses current orchestrators and multikernel management tools.

\subsection{Architectures with multiple kernels} \label{sec:rel-multikernel}

The idea of structuring the OS as a set of independent kernels emerged a few decades ago.
Hive~\cite{chapin1995hive} organizes the OS as an internal distributed system of kernels, called cells,
where each one is responsible for a fixed set of nodes and operates as an independent OS, while they cooperate
to offer applications a single, transparent system image.
Communication between cells occurs through \textit{remote procedure calls} (RPC) and by directly reading the memory of other cells.

Instead of rewriting the OS entirely, Disco~\cite{bugnion1997disco} inserts a small virtual machine monitor
that runs several copies of an OS on the same machine, communicating through common distributed protocols, such as NFS and
TCP/IP, which makes it possible to run specialized operating systems side by side with commercial systems.
Cellular Disco~\cite{govil2000cellular} extends the proposal by dividing the monitor itself into smaller parts, in order to
contain hardware faults, turning the multiprocessor into what would be a virtual cluster.

Corey~\cite{boydwickizer2008corey}, by contrast, starts from the principle that applications should control sharing.
By default, each kernel data structure is used by a single core, and abstractions such as \textit{address ranges},
\textit{kernel cores}, and \textit{shares} allow the application to explicitly indicate what should be shared.
Although it runs a single kernel, Corey was an important foundation for later works, such as Barrelfish.
FOS~\cite{wentzlaff2009fos}, in turn, proposes factoring OS services into fleets of servers spatially distributed
across the cores, with communication exclusively through messages, motivated by the finding that the
Linux physical page allocator did not scale beyond eight cores.
Message-based communication became established in subsequent proposals, including the multikernel implementation based on
Linux~\cite{wang2025multikernel}.

The modern formulation of the concept is presented by Baumann et al.~\cite{baumann2009multikernel}, who propose treating
the machine as a network of independent cores.
The resulting system, Barrelfish, runs one kernel per core and
follows three principles: communication between cores is always explicit, the structure is hardware-neutral, and
state is replicated instead of shared, being kept consistent by monitor processes.

More recently, the proposal for supporting multiple Linux instances on the same machine~\cite{wang2025multikernel}
loads the instances via \textit{kexec}, assigning each one cores and memory that can be reallocated at
run time through a \textit{device tree} and the \textit{hotplug} mechanism~\cite{wang2025lpc}.
This is the implementation of the architecture planned for K-ORCH.

\subsection{Load balancing, observability, isolation, and fundamental concepts} \label{sec:rel-fundamentos}

Deciding where to run each workload is a problem that has long been studied in distributed systems.
Casavant and Kuhl~\cite{casavant1988taxonomy} propose a hierarchical scheduling taxonomy that distinguishes
static from dynamic approaches and classifies the latter according to the distribution of responsibility for decisions and the
cooperation between nodes.
Shivaratri et al.~\cite{shivaratri1992load} decompose load distribution algorithms
into four policies: transfer (when to move a task), selection (which task to move), location
(where to move it), and information (when to collect the state of the nodes).
Mitzenmacher~\cite{mitzenmacher2001power}
shows that choosing the least loaded option among only two random destinations exponentially reduces the expected maximum
load compared to a purely random choice, a result that underpins simple and scalable policies.

In systems with multiple kernels, it was necessary from the start to decide how to distribute the machine's resources
among the instances.
In Hive, the sharing mechanisms are implemented through cooperation between cells,
but the policy is assigned to a user process, the only one with a global view of the system, responsible for deciding
from which cells to allocate memory and how to schedule processes~\cite{chapin1995hive}.
In Disco, the monitor applies
global policies over the virtual machines, moving memory between them and dynamically scheduling the vCPUs
to balance the load~\cite{bugnion1997disco}.
Cellular Disco combines an idle balancer, which uses only
local information, with a periodic balancer, which queries a global load tree~\cite{govil2000cellular}.
In Barrelfish, a knowledge base gathers hardware information and run-time measurements, answering
optimization queries, such as the allocation of drivers to cores~\cite{baumann2009multikernel,schupbach2008diversity}.
In HPC-focused multikernels, management tends to be exposed as an interface.
In IHK, the master kernel takes
control of the physical memory mapping and offers operations to reserve cores and memory, and to initialize,
query, shut down, and destroy instances, but it does not define policies about when to use them~\cite{shimosawa2014ihk}.

Placement decisions depend on information about the state of the nodes.
Mansouri-Samani and
Sloman~\cite{mansourisamani1993monitoring} define a functional model for monitoring distributed systems
divided into four activities (generation, processing, dissemination, and presentation of information) and point to
information staleness as the main problem of monitoring, which directly affects an orchestrator that
makes decisions based on remote state.
In microservices, observability is organized today around three sources:
metrics, logs, and traces~\cite{faseeha2025observability}.
LDMS~\cite{agelastos2014ldms} shows that it is
possible to continuously monitor an entire system, even in HPC, with collector processes on the nodes and
aggregator processes that gather the data, with low overhead on the applications.
Systems with multiple kernels themselves
also have monitoring mechanisms: Barrelfish's knowledge base incorporates measurements such as the
communication latency between pairs of cores~\cite{baumann2009multikernel}, Cellular Disco's periodic balancer depends
on a constantly updated load tree~\cite{govil2000cellular}, and PIKA measures the time of each application
to guide its choices~\cite{beckmann2014pika}.
However, this information is collected and consumed on a single machine,
without sharing among the nodes of a cluster.

A major motivation for dividing the system into multiple kernels is application availability.
In Hive, a hardware or software failure affects only the processes that use the resources of the affected cell.
The system protects each page against writes from other cells, preemptively discards the pages that a failed cell
could have corrupted, and detects failures through heuristics, whose confirmation requires consensus among the surviving
cells~\cite{chapin1995hive}.
Disco treats each virtual machine as a unit of software fault
containment~\cite{bugnion1997disco}, and Cellular Disco extends this containment to hardware faults: in hundreds of
fault-injection experiments, recovery was successful in all cases, in less than half a
second~\cite{govil2000cellular}.
In addition to containing faults, division into multiple kernels is also used to isolate
application performance.
Pisces~\cite{ouyang2015pisces} allows the coexistence, on the same node, of multiple
independent OS environments, called enclaves, each running a lightweight co-kernel that manages a partition
of hardware resources, treating performance isolation as a primary requirement.

\subsection{Orchestration} \label{sec:rel-orquestracao}

Current orchestrators have been enormously influenced by Google's experience with Borg, which was
consolidated in Kubernetes into a declarative model, that is, the user describes the desired state, and control loops
continuously reconcile the observed state with it~\cite{burns2016borg}.
A review of orchestrators such as Kubernetes, Docker Swarm, and Apache Mesos identifies the scarcity of proposals for
resource-utilization-driven orchestration and compares the suitability of each tool~\cite{swetha2025resource}.
Along the same lines, a two-stage container deployment distributes containers among virtual machines with
coarse granularity through a greedy algorithm, while a genetic algorithm divides the resources within
each virtual machine~\cite{lu2025optimized}, a structure similar to choosing a node and then performing local partitioning,
something that will be adopted in this work.

For workloads that require stronger isolation than that of containers, the predominant approach is to integrate lightweight virtual
machines into existing orchestrators. Firecracker~\cite{agache2020firecracker} is a minimal virtual machine
monitor, developed to run functions and containers from different customers in microVMs, and
Kata Containers~\cite{katacontainers} runs each Kubernetes pod in a lightweight virtual machine, integrated
through the container runtime interface (CRI).
Orchestrated microVMs are, therefore, the comparison baseline adopted in this work.

In the multikernel context, instance management remains local. In the Hobbes project, a node manager is
responsible for the lifecycle of distinct enclaves on the same node and for the shared memory mechanisms between
them~\cite{kocoloski2015composition}. For the Linux multikernel proposal, the kerf tool~\cite{multikernel2025kerf}
manages the creation, configuration, monitoring, and destruction of instances, being the tool closest
to an orchestrator, but limited to a single machine.

\section{Methodology} \label{sec:methods}

The methodology used in this work is theoretical-practical in nature, combining analysis of orchestration algorithms and
mechanisms with their respective implementations and an experimental evaluation.
The approach is quantitative, based on identifying the benefits and comparing performance
between applications using measurements of the behavior of the orchestrator and of instances under controlled workloads.
Finally, the research level is exploratory, testing and verifying different possibilities that
can be used to manage this new type of technology, given that distributed orchestration of the multikernel
architecture has not yet been documented in the literature.

% Define multikernel instance
\subsection{Characterize and define the orchestrable unit}\label{subsec:caracterizar-e-definir-a-unidade-orquestravel}

In the first stage, the state of the art will be surveyed and the orchestrable unit will be defined.
The multikernel platform will be studied experimentally regarding boot time, the way cores and
memory are reserved, the instances' access to the network, and possible signals available for detecting failures.
Based on this characterization, what makes up an instance (dedicated
cores, memory region, image, and kernel configuration) and the desired-state model maintained by the orchestrator will be formally defined.

% Implement orchestration prototype with the parent node and agents
\subsection{Design and implementation of the prototype}\label{subsec:projeto-e-implementacao-do-prototipo}

The second stage accomplishes the objectives of designing the architecture and implementing the lifecycle and placement.
The control plane will be implemented in order to receive instance requests, described by state.
The agent, running on the host kernel of each node, will be commanded by the control plane, performing resource partitioning and the initialization, stopping, and removal of
instances, reporting the state of each of them.
Initially, two basic placement policies will be implemented: \textit{first-fit} and choosing the less
loaded node between two randomly drawn ones~\cite{mitzenmacher2001power}.

% evolve features to integrate observability, balancing, and others
\subsection{Resource partitioning, monitoring, and recovery}\label{subsec:particionamento-de-recursos-monitoramento-e-recuperacao}

In the third stage, resource partitioning and monitoring will be carried out.
The orchestrator will operate in two modes: static, with predefined partitions on each node, and dynamic, with cores and memory partitioned
on each request.
The agent will be a periodic collector, covering the state and resource usage of each
instance and of the nodes; the control plane must aggregate this information and make decisions based on it.
Failures will be detected by the absence of signals from the agent and by events; recovery will consist of restarting the instance
on the same node or relocating it to another one that is available and healthy.

% perform quantitative and comparative analysis of the tool
\subsection{Analysis and refinements}\label{subsec:analysis-and-refinements}

The fourth stage consists of carrying out the experimental evaluation, which will be conducted on a set of physical nodes connected on
the same local network, in addition to testing the tool with specific experiments:
\begin{itemize}
  \item \textbf{Research and explore workloads:} Workloads will be selected from established benchmarks, according to
  criteria fixed before the tests are executed, and must cover both scenarios in which the multikernel architecture tends to
  be favored and scenarios in which it tends to be disfavored.
  In this way, the results are not conditioned on a choice of workloads that favors the proposal and can provide clear and meaningful results.
  \item \textbf{Provisioning and orchestrator cost:} time between the request and the instance's response,
  broken down into partitioning, kernel boot, and availability, compared to the boot
  time of microVMs. Also analyze CPU and memory consumption of the agent and the control plane.
  \item \textbf{Density and interference:} run latency-sensitive, CPU-bound, and memory-bound workloads alongside other
  applications, comparing multikernel and microVMs.
  \item \textbf{Static and dynamic partitioning:} sequences of requests with many different sizes,
  measuring the utilization of cores and memory on each node, rejected requests, and provisioning time.
\end{itemize}

With this data collected, an analysis will be performed and more specific modifications to the orchestrator may be made,
in order to improve its performance.

\section{Expected Results} \label{sec:results}

The main expected result is a functional prototype of the K-ORCH tool, capable of provisioning, placing, and monitoring multikernel
instances across multiple nodes.
The proposal seeks to understand and explore the best paths for the specific characteristics of the multikernel architecture, creating
this robust tool for the orchestration of distributed systems.

It is also intended to create a record of the parameters used and of the greatest difficulties in carrying out orchestration,
mapping and documenting algorithms, main ideas, performance, and other metrics.
In this way, it will also be possible to offer a characterization of the costs and limits of distributed multikernel orchestration.
Finally, the intention is to disseminate the proposal and the results obtained through a paper submitted
to a relevant conference or journal in this research area.

\section{Schedule} \label{sec:crono}

Table \ref{tab:crono} presents the schedule organized into 8 quarters.
Quarters one and two focus on studying the multikernel platform and formally defining the orchestrable unit used, creating
a reproducible experimental environment.
Quarters three and four will be used to design the general operation of the orchestrator, and then implement the orchestration
agents in order to carry out the application lifecycle and placement, with a functional orchestration unit as a milestone
at the end of the year, serving as a proof of concept.
Quarters five and six will be dedicated to resource partitioning, monitoring, and
instance recovery, as well as preparing the environment for comparison with microVMs.
Finally, in quarters seven and eight, tests and analyses will be carried out regarding the performance of the orchestration
and of the chosen applications, that is, an experimental evaluation focused on refining the orchestrator and consolidating the results.

\begin{table}[ht]
  \centering
  \caption{Task schedule organized by quarter.}
  \label{tab:crono}
  \small
  \renewcommand{\arraystretch}{1.25}
  \begin{tabular}{p{7.2cm}cccccccc}
    \hline
    \textbf{Activity} & Q1 & Q2 & Q3 & Q4 & Q5 & Q6 & Q7 & Q8 \\
    \hline
    Survey of the state of the art                                   & $\bullet$ & $\bullet$ &           &           &           &           &           &           \\
    Experimental characterization of the multikernel platform        & $\bullet$ & $\bullet$ &           &           &           &           &           &           \\
    Definition of the orchestrable unit and state model              &           & $\bullet$ &           &           &           &           &           &           \\
    Architecture design                                              &           &           & $\bullet$ &           &           &           &           &           \\
    Implementation of lifecycle and placement                        &           &           & $\bullet$ & $\bullet$ &           &           &           &           \\
    Implementation of resource partitioning                          &           &           &           &           & $\bullet$ & $\bullet$ &           &           \\
    Implementation of instance monitoring and recovery               &           &           &           &           & $\bullet$ & $\bullet$ &           &           \\
    Preparation of the evaluation environment and reference microVMs &           &           &           &           & $\bullet$ & $\bullet$ &           &           \\
    Experimental and comparative evaluation                          &           &           &           &           &           &           & $\bullet$ & $\bullet$ \\
    Orchestrator refinement                                          &           &           &           &           &           &           & $\bullet$ & $\bullet$ \\
    Writing of papers and the final text                             &           &           &           & $\bullet$ &           & $\bullet$ & $\bullet$ & $\bullet$ \\
    \hline
  \end{tabular}
\end{table}

\section{Final Remarks} \label{sec:cons-final}

This work proposes K-ORCH, an orchestrator that extends the management of multikernel instances, currently restricted to a
single machine, to environments with multiple nodes.
The proposal combines the model of current orchestrators with
the partitioning mechanisms specific to the multikernel architecture and will be evaluated against orchestrated microVMs in terms of
provisioning cost, density, and performance interference between instances.
The expected contribution is to provide a third way for environments that need the best of containers and
virtual machines, providing enough abstraction to scale this architecture to large production systems
with ease and reliability.
In addition, it is expected to characterize the costs and limits of distributed orchestration of the multikernel architecture.

\bibliographystyle{sbc}

\bibliography{sbc-template}

\end{document}
